\documentclass[11pt]{article}
\usepackage[letterpaper,margin=1in]{geometry}
\usepackage[T1]{fontenc}
\usepackage{lmodern,microtype}
\usepackage{amsmath,amssymb,amsthm,mathtools}
\usepackage{booktabs,longtable,array}
\usepackage{xcolor,enumitem,fancyhdr}
\usepackage[hyphens]{url}
\usepackage[colorlinks=true,linkcolor=blue!45!black,citecolor=blue!45!black,urlcolor=blue!45!black]{hyperref}
\usepackage{aliascnt}
\usepackage[nameinlink,noabbrev]{cleveref}
\numberwithin{equation}{section}
\newtheorem{theorem}{Theorem}[section]
\newaliascnt{lemma}{theorem}
\newtheorem{lemma}[lemma]{Lemma}
\aliascntresetthe{lemma}
\crefname{lemma}{lemma}{lemmas}
\Crefname{lemma}{Lemma}{Lemmas}
\newaliascnt{proposition}{theorem}
\newtheorem{proposition}[proposition]{Proposition}
\aliascntresetthe{proposition}
\crefname{proposition}{proposition}{propositions}
\Crefname{proposition}{Proposition}{Propositions}
\newaliascnt{corollary}{theorem}
\newtheorem{corollary}[corollary]{Corollary}
\aliascntresetthe{corollary}
\crefname{corollary}{corollary}{corollarys}
\Crefname{corollary}{Corollary}{Corollarys}
\newaliascnt{conjecture}{theorem}
\newtheorem{conjecture}[conjecture]{Conjecture}
\aliascntresetthe{conjecture}
\crefname{conjecture}{conjecture}{conjectures}
\Crefname{conjecture}{Conjecture}{Conjectures}
\theoremstyle{definition}
\newaliascnt{definition}{theorem}
\newtheorem{definition}[definition]{Definition}
\aliascntresetthe{definition}
\crefname{definition}{definition}{definitions}
\Crefname{definition}{Definition}{Definitions}
\newaliascnt{remark}{theorem}
\newtheorem{remark}[remark]{Remark}
\aliascntresetthe{remark}
\crefname{remark}{remark}{remarks}
\Crefname{remark}{Remark}{Remarks}
\newcommand{\Li}{\operatorname{Li}}
\newcommand{\dd}{\,\mathrm d}
\newcommand{\R}{\mathbb R}
\newcommand{\N}{\mathbb N}
\newcommand{\D}{\mathcal D}
\newcommand{\lc}{\rho_{\mathrm c}}
\newcommand{\rs}{\rho_*}
\newcommand{\ac}{a_{\mathrm c}}
\newcommand{\aL}{a_{\mathrm L}}
\newcommand{\aM}{a_{\mathrm M}}
\newcommand{\aO}{a_{\mathrm O}}
\newcommand{\ordli}[2]{\Li_{#1}^{\langle #2\rangle}}
\newcommand{\sha}{\texttt{cc34f73596336f2466d9754cb0f3635bd2bedade}}
\setlength{\headheight}{14pt}
\pagestyle{fancy}
\fancyhf{}
\fancyhead[L]{\small Lerch zeros: global exclusion and unit-span rigidity}
\fancyhead[R]{\small ProveIt continuation}
\fancyfoot[C]{\thepage}
\setlength{\emergencystretch}{2em}
\setlist{itemsep=3pt,topsep=5pt}
\hypersetup{pdftitle={A Complete Index-Three Lerch Phase Diagram},pdfauthor={OpenAI-assisted research continuation for Vladimir Reshetnikov}}
\title{A Complete Index-Three Lerch Phase Diagram\\[5pt]
\large Global exclusion, sharp exterior barriers, and polylogarithmic unit-span identities}
\author{Research continuation for Vladimir Reshetnikov's \texttt{ProveIt} project\\
\small OpenAI-assisted mathematical research draft}
\date{October 9, 2026\quad (America/Los\_Angeles)}
\begin{document}
\maketitle
\thispagestyle{empty}
\begin{abstract}
We resolve the remaining global-exclusion conjecture in the earlier report
\emph{Lerch Endpoint Singularities and Stieltjes Zero Bifurcations}.
For
\[
 F(a,\rho)=\sum_{m=0}^{\infty}\rho^m
 \frac{\log^2(a+m)\,[3-\log(a+m)]}{(a+m)^2},
 \qquad a>0,\quad 0\leq\rho\leq1,
\]
there is exactly one outer positive zero for every deformation parameter.
Consequently the previously localized index-three fold is the complete global
phase transition: one simple zero below the unique critical parameter, a double
zero and a simple zero at the critical parameter, and three simple zeros above it.
The proof combines an analytic exterior differential inequality with an exact
rational interval certificate on one compact interval, uniformly over an entire
parameter interval. It is not a numerical scan of a two-dimensional region.

A general theorem supplies exterior one-crossing barriers at every Laurent
index, with a sharp cutoff for the fixed differential operator
$\partial_a+e^{-n}$ and an exponentially accurate asymptotic description of that
cutoff. A separate rigidity theorem shows that the three-zero span is at most
one. Equality occurs at a unique parameter where two zeros are exactly $1$ and
$2$; the span has a nondegenerate global maximum there. The resonance and its
curvature admit explicit formulas in polylogarithmic order derivatives.
We also prove an all-odd-index restriction on integer-separated zeros and a
family of strict polylogarithmic tail inequalities. All proof-bearing finite
computations use integer interval arithmetic and a proved Euler--Maclaurin
remainder; independent high-precision diagnostics are labeled separately.
\end{abstract}
\medskip
\noindent\textbf{Status.} Ordinary mathematical proofs plus reproducible exact-arithmetic
certificates; not proof-assistant verified or externally refereed. Novelty is
asserted relative to the inspected manuscript and named companion, not as an
exhaustive priority claim. No remote repository files were modified.
\newpage
\tableofcontents
\newpage

\section{Scope, provenance, and the resolved question}
\label{sec:scope}
The source is Vladimir Reshetnikov's \texttt{ProveIt} repository, with manuscript
subtree
\begin{center}\small
\path{Analysis/Polylogarithms/docs/manuscript}.
\end{center}
The inspected branch was pinned to commit
\begin{center}\small\sha.\end{center}
GitHub records its UTC commit time as October 10, 2026, 02:47:46, which is
October 9, 2026 in the user's Pacific timezone.
The zero-geometry source chapter has Git blob identifier
\texttt{7f41eb87ed5e0bf90a2dffa489a8e2cd552056be}.
We inspected that chapter, relevant passages of \texttt{10-discovery.tex}, the
manuscript directory inventory, and the earlier companion report
\cite{repo,previous}. The latter was retrieved as its LaTeX source; its hash is
recorded in the package. A same-named report archive is listed in
\texttt{docs/incoming} at the pinned commit, but its archive bytes were not used.

The manuscript already proves the positive-kernel representation, the
$n$-zero upper bound, eventual saturation, and the low-index classical counts.
The companion already proves a unique \emph{left-hand} index-three fold, but
explicitly leaves open the possible existence of an additional outer pair below
that fold. Its Conjecture 12.1 asks for the complete global count.
This distinction matters: nested negative sets and a local fold do not,
without further information, exclude a disconnected pair of zeros elsewhere.

The principal new deductions are organized as follows.
\begin{center}
\small
\begin{tabular}{p{.27\linewidth}p{.62\linewidth}}
\toprule
Result & Content and evidence\\
\midrule
Global outer exclusion & Exactly one simple zero in $(9/5,e^3)$ for each
$0<\rho\leq1$; analytic barrier plus a finite exact certificate.\\
Complete phase diagram & Resolves the companion's Conjecture 12.1, including
multiplicity and branch directions.\\
All-index exterior theorem & A strict first-order differential inequality,
a sharp cutoff for its fixed coefficient, and a cutoff asymptotic.\\
Unit-span rigidity & The span of the three zeros is at most $1$, with a unique,
nondegenerate equality case at roots $1$ and $2$.\\
Polylogarithmic identities & A unique order-derivative resonance, exact common
velocity and curvature formulas, and strict tail-moment inequalities.\\
Integer-gap restriction & For every odd index at derivative order one, the
only possible positive integer gap between zeros is the pair $1,2$.\\
\bottomrule
\end{tabular}
\end{center}

The recent $S_4$ proof is not repeated here. The weight-seven $S_6$ reduction,
all-index saturation thresholds, and a classification of higher-index folds
are not claimed to be settled. The classical Lerch shift equation, Mellin
representation, gamma product, and Euler--Maclaurin formula are background
rather than new discoveries \cite{dlmf-lerch,dlmf-hurwitz,dlmf-gamma,dlmf-em,coffey}.

\section{The derivative family and conventions}
\label{sec:family}
We use the Laurent convention
\begin{equation}
 \zeta(s,a)=\frac1{s-1}+\sum_{n\geq0}\frac{(-1)^n}{n!}\gamma_n(a)(s-1)^n.
\end{equation}
For $0\leq\rho<1$, define
\[
 \ell_n(\rho,a)=\sum_{m\geq0}\rho^m\frac{\log^n(a+m)}{a+m},
 \qquad
 \Phi(\rho,s,a)=\sum_{m\geq0}\frac{\rho^m}{(a+m)^s}.
\]
Only the differentiated family extends by this series to $\rho=1$.
Writing
\begin{equation}\label{eq:gn}
 g_n(x)=\frac{\log^{n-1}x\,(n-\log x)}{x^2},\qquad
 U_n(a,\rho)=\sum_{m\geq0}\rho^m g_n(a+m),
\end{equation}
we have $U_n=\partial_a\ell_n$ for $\rho<1$, and
$U_n(a,1)=\gamma_n'(a)$. A prime on $\gamma_n$ differentiates its parameter
$a$. In contrast, $\zeta^{(j)}(s,a)$ below differentiates the complex order $s$.
The identity at $\rho=1$ follows either from the standard parameter identity
$\partial_a\zeta(s,a)=-s\zeta(s+1,a)$ or directly from its Laurent expansion.

Our main function is
\begin{equation}\label{eq:Fdef}
 F(a,\rho)=U_3(a,\rho),\qquad
 f(x)=g_3(x)=x^{-2}\log^2x(3-\log x).
\end{equation}
It agrees with the manuscript's $F_{3,1}^{\rho}$, since the sign normalization
$(-1)^{n+k}$ is positive at $(n,k)=(3,1)$.
For fixed positive $a$, and uniformly on compact positive $a$-intervals,
the series defining $F$ and all its fixed-order $a$-derivatives converge
absolutely, uniformly in $0\leq\rho\leq1$. This follows from the majorants
$O(m^{-2-r}(1+\log^3m))$ after $r$ parameter derivatives.
Thus $F$ is continuous up to $\rho=1$ and is real analytic in $(a,\rho)$ for
$a>0$, $0<\rho<1$. No finite $\rho$-derivative at $\rho=1$ is assumed.

Several elementary facts will be used repeatedly:
\begin{align}
 f(x)&\geq0 &&(0<x\leq e^3),\label{eq:signf}\\
 f(x)&<0 &&(x>e^3),\notag\\
 f(x)=0&\Longleftrightarrow x\in\{1,e^3\},\notag\\
 f'(x)&=x^{-3}\bigl(2\log^3x-9\log^2x+6\log x\bigr),\label{eq:df}\\
 F(a+1,\rho)&=\frac{F(a,\rho)-f(a)}{\rho}
 &&(0<\rho\leq1).\label{eq:shift}
\end{align}
The last equation is the classical Lerch shift relation in the present
coordinate. Since the tails with $m\geq1$ are uniformly bounded for $0<a\leq1$,
$f(a)\to+\infty$ proves that $F(a,\rho)>0$ at sufficiently small $a$, uniformly
in $\rho$. For $a>e^3$, every summand is negative. For $\rho>0$ the value at
$a=e^3$ is also strictly negative.

\section{Two variation principles and exact anchors}
\label{sec:variation}
We reproduce the ingredients needed for a self-contained zero count.

\subsection{The multiplicity bound in the parameter \texorpdfstring{$a$}{a}}
Define the reciprocal-gamma cubic
\begin{equation}\label{eq:Q3}
 Q_3(x)=(x+\gamma)^3-3\zeta(2)(x+\gamma)+2\zeta(3).
\end{equation}
It is the coefficient of $t^3/3!$ in $e^{xt}/\Gamma(1+t)$.
The standard gamma expansion gives this coefficient explicitly
\cite{dlmf-gamma}. The Mellin representation gives
\begin{equation}\label{eq:laplace}
 F(a,\rho)=\int_0^\infty
 \frac{x e^{-ax}}{1-\rho e^{-x}}Q_3(\log x)\dd x.
\end{equation}
Indeed,
\[
 (1+t)\Phi(\rho,2+t,a)
 =\frac1{\Gamma(1+t)}\int_0^\infty
       \frac{x^{1+t}e^{-ax}}{1-\rho e^{-x}}\dd x,
\]
and $3![t^3]$ on the left is exactly the defining series for $F$.
For $\rho=1$, the same equation holds in the convergent Hurwitz half-plane.
Near zero the integrand in \eqref{eq:laplace} is bounded in absolute value by a
constant times $1+|\log x|^3$, uniformly in $\rho$; at infinity the exponential
factor dominates. These bounds justify extraction and all needed derivatives.

\begin{lemma}[Laplace variation with multiplicity]\label{lem:laplacevariation}
If an integrable signed kernel has $N$ sign changes on $(0,\infty)$, with
nonzero constant sign between successive changes, and all required
exponentially weighted moments exist, its nonzero Laplace transform has at
most $N$ positive zeros, counted with multiplicity.
\end{lemma}
\begin{proof}
For no sign changes the transform has a strict sign. If $x_0$ is a sign-change
point, multiplication of the kernel by $x_0-x$ removes that change and no
others. The corresponding transform is
\[
 (\partial_a+x_0)G(a)=e^{-x_0a}\frac{\mathrm d}{\mathrm da}
                         \bigl(e^{x_0a}G(a)\bigr).
\]
Induction bounds its zeros by $N-1$. Any finite collection of zeros of $G$
with total multiplicity $M$ forces at least $M-1$ zeros of the displayed
weighted derivative by Rolle's theorem, including inherited multiplicities.
Therefore $M\leq N$. The same induction proves nonvanishing, since an
identically zero transform would have an identically zero weighted derivative.
\end{proof}

\begin{proposition}[Three-zero upper bound]\label{prop:three}
For every $0\leq\rho\leq1$, $F(\cdot,\rho)$ has at most three positive zeros,
counted with multiplicity.
\end{proposition}
\begin{proof}
The discriminant of the centered cubic in \eqref{eq:Q3} is
$108(\zeta(2)^3-\zeta(3)^2)>0$. For an elementary check of positivity,
$\zeta(2)>5/4$ and
$\zeta(3)<1+1/8+\int_2^\infty x^{-3}\dd x=5/4$ suffice.
Thus $Q_3$ has three distinct real roots. The positive multiplier in
\eqref{eq:laplace} leaves exactly three kernel sign changes, and
\cref{lem:laplacevariation} applies. The positive small-$a$ sign also
excludes an identically zero transform.
\end{proof}

\subsection{One crossing in the deformation parameter}
\begin{lemma}[Geometric-weight crossing]\label{lem:rhocross}
Fix $0<a<e^3$. The function $\rho\mapsto F(a,\rho)$ has at most one zero in
$(0,1)$, and at every such zero
\begin{equation}\label{eq:Frhoneg}
 F_\rho(a,\rho)<0.
\end{equation}
If $F(a,\rho_0)>0$ for some $0<\rho_0\leq1$, then $F(a,\rho)>0$ for
$0<\rho\leq\rho_0$. Negativity, once attained, persists for larger $\rho$.
\end{lemma}
\begin{proof}
The coefficient signs in $\sum\rho^m f(a+m)$ change once from positive to
negative, after zeros are omitted. Choose $\lambda$ strictly between the
last nonnegative block and the first negative block. At a zero,
\[
 \rho F_\rho=\sum_{m\geq0}(m-\lambda)\rho^m f(a+m)<0.
\]
Every nonzero term is negative. Existence of a second zero would require
another crossing direction, or a zero with vanishing derivative, which is
impossible. Alternatively, comparison with a separating power of
$\rho/\rho_0$ proves the sign persistence directly. The first nonzero
coefficient is positive, so the sign near $\rho=0$ is positive, including
when $a=1$ and the constant term vanishes.
\end{proof}

\subsection{Four exact endpoint anchors}
The package independently recomputes the following enclosures. Every finite
decimal here is an exact rational endpoint, with outward rounding:
\begin{center}
\begin{tabular}{cl}
\toprule
$a$ & Enclosure for $F(a,1)$\\
\midrule
$4/5$ & $[0.2560298104220831,\ 0.2560298104220832]$\\
$1$ & $[-0.0323050999463418,\ -0.0323050999463417]$\\
$3/2$ & $[0.0249055835689205,\ 0.0249055835689206]$\\
$9/5$ & $[0.0052639818935007,\ 0.0052639818935008]$\\
\bottomrule
\end{tabular}
\end{center}
The first three signs are consistent with the manuscript and companion;
the fourth is the useful additional separator for the present exclusion.
The proof of these inclusions, rather than their decimal appearance, is
specified in \cref{sec:certificate} and implemented in
\texttt{code/verify\_exact.py}.

\begin{corollary}[Common zero-free regions]\label{cor:separator}
For every $0<\rho\leq1$,
\begin{equation}\label{eq:zerofree}
 F(a,\rho)>0\quad\text{on }(0,4/5]\ \text{and on }[3/2,9/5].
\end{equation}
At $\rho=1$ there are precisely three simple zeros: one in $(4/5,1)$, one
in $(1,3/2)$, and one in $(9/5,e^3)$.
\end{corollary}
\begin{proof}
The four signs and $F(e^3,1)<0$ produce three disjoint sign-changing
intervals. \Cref{prop:three} makes these all the zeros and all simple.
The positive small-$a$ sign and absence of further zeros give positivity on
the stated closed separators at $\rho=1$. \Cref{lem:rhocross} propagates it
to smaller positive $\rho$.
\end{proof}

\section{An all-index exterior differential barrier}
\label{sec:barrier}
This section applies to every integer $n\geq2$, not just to the cubic case.

\begin{theorem}[An explicit exterior one-crossing inequality]\label{thm:barrier}
For $n\geq2$, $x\geq e^{(n-1)/2}$,
\begin{equation}\label{eq:barrierpoint}
 g_n'(x)+e^{-n}g_n(x)<0.
\end{equation}
Consequently, for every $0\leq\rho\leq1$,
\begin{equation}\label{eq:barriersum}
 \partial_a U_n(a,\rho)+e^{-n}U_n(a,\rho)<0,
 \qquad a\geq e^{(n-1)/2}.
\end{equation}
The function $e^{a/e^n}U_n(a,\rho)$ is strictly decreasing there. In
particular, there is at most one zero in this exterior interval, and any
such zero is simple and crosses from positive to negative.
\end{theorem}
\begin{proof}
Set $y=\log x$. Up to the positive factor $x^{-3}y^{n-2}$, the left side of
\eqref{eq:barrierpoint} is
\begin{equation}\label{eq:Rn}
 R_n(y)=2y^2-3ny+n(n-1)+e^{y-n}y(n-y).
\end{equation}
For $y=n+r$, $r\geq0$, using $e^r\geq1+r$ gives
\[
 R_n(n+r)=2r^2+nr-n-e^r(n+r)r
 \leq-r^3-(n-1)r^2-n<0.
\]
For $y=n-t$, $0\leq t\leq(n+1)/2$, the factor $(n-t)t$ is nonnegative.
Since $e^{-t}\leq(1+t)^{-1}$,
\[
 (1+t)R_n(n-t)
 \leq 2t^3+(1-n)t^2-nt-n
 \leq 2t^2-nt-n.
\]
The last quadratic is convex. Its values at the endpoints of
$[0,(n+1)/2]$ are $-n$ and $(1-n)/2$, both negative. Hence it is negative
throughout the interval. This proves the pointwise inequality. Termwise
summation with nonnegative weights and the established convergence proves
\eqref{eq:barriersum}; its integrating factor gives the remaining claims.
\end{proof}

\begin{corollary}[Positive-measure extension]\label{cor:measure}
Let $\nu$ be a nonzero nonnegative measure on $[0,\infty)$, and suppose
\[
 H_n(a)=\int_{[0,\infty)}g_n(a+t)\dd\nu(t)
\]
and its first derivative converge absolutely, locally uniformly on the
region in question. Then $H_n'+e^{-n}H_n<0$ for
$a\geq e^{(n-1)/2}$. The same exterior one-crossing conclusion holds.
\end{corollary}
\begin{proof}
Integrate the strict pointwise inequality. A strictly negative integrable
function has strictly negative integral against a nonzero positive measure.
\end{proof}

The positivity of the weights is essential. This argument does not extend
formally to nonreal roots of unity or to arbitrary signed colors.

\subsection{The sharp cutoff for the fixed coefficient \texorpdfstring{$e^{-n}$}{exp(-n)}}
The convenient exponential cutoff above can be improved substantially.
Define
\[
 \Delta_n=\sqrt{n^2+8n},\qquad
 \alpha_n=\frac{3n-\Delta_n}{4}.
\]
The number $\alpha_n$ is the smaller positive root of the polynomial part
of $R_n$.

\begin{theorem}[Sharp fixed-operator cutoff]\label{thm:sharp}
For every $n\geq2$, there is a unique
\begin{equation}\label{eq:xirange}
 \xi_n\in(\alpha_n,(n-1)/2)
\end{equation}
with $R_n(\xi_n)=0$. One has $R_n(y)<0$ for every $y>\xi_n$ and
$R_n(y)>0$ for $0<y<\xi_n$. Thus
\[
 g_n'(x)+e^{-n}g_n(x)<0\qquad(x>e^{\xi_n}),
\]
and $e^{\xi_n}$ is the sharp lower cutoff for this strict pointwise inequality.
At $x=e^{\xi_n}$ the left side is zero. For geometric weights, the summed
inequality is strict even at $a=e^{\xi_n}$ when $\rho>0$; at $\rho=0$ that
endpoint has equality in the differential inequality, not a zero of $U_n$.
\end{theorem}
\begin{proof}
Write $H_n(t)=R_n(n-t)$. The preceding proof shows $H_n(t)<0$ for
$0\leq t\leq(n+1)/2$. On $[(n+1)/2,n]$,
\[
 H_n''(t)=4+e^{-t}\bigl(-t^2+(n+4)t-2n-2\bigr)>4.
\]
The quadratic in parentheses is concave, with endpoint values
$(n^2-1)/4$ and $2(n-1)$, so it is strictly positive. Since
$H_n(n)=n(n-1)>0$, strict convexity gives precisely one zero in this
interval: after a first crossing from a negative value, its derivative is
positive and remains so. This determines a unique $\xi_n$ in
$(0,(n-1)/2)$ and gives the asserted signs below $y=n$; the proof of
\cref{thm:barrier} supplies the sign above $n$.
Finally $R_n(\alpha_n)=e^{\alpha_n-n}\alpha_n(n-\alpha_n)>0$, proving
$\alpha_n<\xi_n$. Below this cutoff the pointwise inequality fails, so no
smaller cutoff works for all positive measures: take a point mass at zero.
For $\rho>0$, the summand at $m=1$ is strictly negative even when the $m=0$
summand has equality, proving the endpoint assertion for the sum.
\end{proof}

\begin{theorem}[Asymptotic location of the sharp cutoff]\label{thm:xiasympt}
As $n\to\infty$,
\begin{equation}\label{eq:xiasympt}
 \xi_n=\alpha_n+
 \frac{\alpha_n(n-\alpha_n)}{\Delta_n}e^{\alpha_n-n}
 +O(n^2e^{-n}).
\end{equation}
In particular, the difference between the sharp cutoff in logarithmic
coordinates and the first stationary point of $g_n$ is exponentially small.
\end{theorem}
\begin{proof}
Let $\xi_n=\alpha_n+\delta$ and $E=e^{\alpha_n-n}$. The exact equation is
\begin{equation}\label{eq:deltaeq}
 0=-\Delta_n\delta+2\delta^2+
 E e^\delta\bigl(\alpha_n(n-\alpha_n)+(n-2\alpha_n)\delta-\delta^2\bigr).
\end{equation}
Here $\Delta_n\asymp n$, $\alpha_n=n/2-1+O(n^{-1})$, and
$E\asymp e^{-n/2}$. Evaluating the right side first at zero and then at
twice $E\alpha_n(n-\alpha_n)/\Delta_n$ gives opposite signs for all
sufficiently large $n$: the omitted terms have size $O(n^3E^2)$ whereas
the main magnitude is of order $n^2E$. Uniqueness from \cref{thm:sharp}
therefore gives $\delta=O(nE)$. Substitute that estimate back into
\eqref{eq:deltaeq}. After division by $\Delta_n$, all terms besides the
first correction contribute $O(n^2E^2)=O(n^2e^{-n})$, as claimed.
\end{proof}

For orientation only, the non-certified diagnostics give
$\xi_2\approx0.4122953342$, $\xi_3\approx0.8515049859$, and
$e^{\xi_3}\approx2.3431706385$. The proof of the cubic phase diagram below
uses the simpler cutoff $e$, so none of these decimal approximations is
proof-bearing. Sharpness here concerns the fixed operator
$\partial_a+e^{-n}$, not an optimization over all positive operator coefficients.

\section{A compact-interval exclusion with exact arithmetic}
\label{sec:compact}
We now bridge the gap between the positive separator $9/5$ and the analytic
exterior region $a\geq e$.

\begin{lemma}[An elementary low-$\rho$ positivity bound]\label{lem:lowrho}
For $0\leq\rho\leq9/10$ and $9/5\leq a\leq11/4$,
\begin{equation}\label{eq:lowrhomargin}
 F(a,\rho)\geq \frac6{121}-\frac{33}{160}\left(\frac9{10}\right)^{18}>0.
\end{equation}
For $4/5\leq a\leq3/2$ and $0<\rho\leq9/10$, the same positive lower
bound multiplied by $\rho$ is valid.
\end{lemma}
\begin{proof}
The elementary exponential series proves $e<11/4<3$ and
$e^3>(163/60)^3>20$. For $x\in[9/5,11/4]$, these inequalities imply
$1/2<\log x<3/2$, and consequently
\[
 f(x)\geq\frac{(1/2)^2(3/2)}{(11/4)^2}=\frac6{121}.
\]
For $x=e^{3+t}$, $t\geq0$, the magnitude of the negative part is
\[
 |f(x)|=e^{-6}(t^3+6t^2+9t)e^{-2t}
 \leq\frac{33}{4}e^{-6}<\frac{33}{1600}.
\]
We used $t^j e^{-2t}\leq j!/2^j$, which follows from the exponential
series, separately for $j=1,2,3$.
For $a\leq11/4$, no negative summand occurs before $m=18$.
Keep the $m=0$ term, discard all other positive terms, and bound the
negative geometric tail to obtain \eqref{eq:lowrhomargin}.

For the second interval use the $m=1$ term instead: $a+1\in[9/5,5/2]$.
The $m=0$ term is nonnegative, and no negative term occurs before $m=19$.
Factoring out one power of $\rho$ gives the stated estimate. The exact
positive margin is
\[
 \frac{360672121259082509847}{19360000000000000000000}.
\]
\end{proof}

\begin{lemma}[Certified derivative strip]\label{lem:strip}
For every $9/5\leq a\leq11/4$ and $9/10\leq\rho\leq1$,
\begin{equation}\label{eq:strip}
 \partial_aF(a,\rho)<-\frac1{40}<0.
\end{equation}
\end{lemma}
\begin{proof}
Partition $[9/5,11/4]$ into $128$ equal rational closed intervals.
For each such interval $I$, compute an outward enclosure of
$f'(I+m)$ for $0\leq m<64$ using \eqref{eq:df} and rational logarithm
bounds. Denote its upper endpoint by $u_m(I)$. For the entire parameter
interval, an upper bound for the weighted summand is
\[
 \rho^m f'(a+m)\leq
 \begin{cases}(9/10)^m u_m(I),&u_m(I)\leq0,\\
 u_m(I),&u_m(I)>0.
 \end{cases}
\]
This remains valid if the underlying interval straddles zero.
For $b=64+9/5=329/5$, the uniform tail is bounded above by
\begin{equation}\label{eq:derivtail}
 T_b=\frac{2\log^3b}{b^3}
 +\frac{\log^3b+\frac32\log^2b+\frac32\log b+\frac34}{b^2}.
\end{equation}
Indeed, $f'(x)\leq2\log^3x/x^3$ for $x\geq b$; the latter positive
majorant is decreasing, so a first term plus its integral bounds the
sum. Negative tail terms, if any, only improve this upper bound.

Every rational interval computation is specified in \cref{sec:certificate}
and recomputed by the accompanying standard-library verifier. The largest
upper bound among the $128$ cells, including $T_b$, is enclosed by
\[
 [-0.0256107476350181,\,-0.0256107476350180],
\]
which is strictly below $-1/40$. The tail bound alone is enclosed by
\[ [0.0251601160205428,0.0251601160205429]. \]
The JSON certificate records the rational endpoints and upper bound of
\emph{every} cell. Thus the finite verification covers the complete closed
strip, not a set of sampled points.
\end{proof}

\begin{remark}
The certificate is deliberately elementary rather than optimized. It uses
only $64$ summands and a rough tail majorant. The proof does not need a
high-precision evaluation of $\rho_{\mathrm c}$, a near-endpoint Lerch
summation, or an interval treatment of the singular $\rho$-derivative at one.
\end{remark}

\section{Unique outer zero and the complete phase diagram}
\label{sec:phase}

\begin{theorem}[Global outer exclusion]\label{thm:outer}
For every $0<\rho\leq1$, $F(\cdot,\rho)$ has exactly one zero in
$(9/5,\infty)$. It lies in $(9/5,e^3)$, is simple, and has negative
$a$-derivative. Denote it by $\aO(\rho)$.
It extends continuously to $\aO(0)=e^3$, is analytic for $0\leq\rho<1$,
and is strictly decreasing on $[0,1]$.
\end{theorem}
\begin{proof}
Existence between $9/5$ and $e^3$ follows from \cref{cor:separator} and
\eqref{eq:signf}. If $0<\rho\leq9/10$, \cref{lem:lowrho} excludes zeros
up to $11/4>e$; \cref{thm:barrier} with $n=3$ allows at most one zero
thereafter, with negative derivative.

If $9/10\leq\rho\leq1$, \cref{lem:strip} makes $F$ strictly decreasing
on $[9/5,e]$. Above $e$, the function $e^{a/e^3}F(a,\rho)$ is strictly
decreasing. A crossing before $e$ therefore leaves a negative value
which cannot cross back above $e$; if no crossing occurs before $e$,
the exterior barrier permits at most one afterwards. A crossing exactly
at $e$ also has negative derivative. This proves uniqueness and simplicity
in every case.

For $0<\rho<1$, \cref{lem:rhocross} gives $F_\rho<0$ at the zero,
whereas $F_a<0$. Implicit differentiation gives
\begin{equation}\label{eq:outervelocity}
 \aO'(\rho)=-\frac{F_\rho(\aO(\rho),\rho)}{F_a(\aO(\rho),\rho)}<0.
\end{equation}
At $\rho=0$, the outer elementary zero $e^3$ is simple, so analytic
continuation is available there. Uniform convergence, the compact
bracket $[9/5,e^3]$, and uniqueness imply continuity at both endpoints.
Strict decrease on the open interval extends to strict decrease on the
closed interval by continuity and the intermediate comparisons.
\end{proof}

\begin{theorem}[Complete index-three phase diagram]\label{thm:phase}
There is a unique $\lc\in(9/10,1)$ and a unique $\ac\in(4/5,3/2)$ such that
\[
 F(\ac,\lc)=F_a(\ac,\lc)=0.
\]
The total positive-zero count is as follows:
\begin{center}
\begin{tabular}{cll}
\toprule
Parameter & Distinct zeros & Multiplicities\\
\midrule
$\rho=0$ & $1,e^3$ & $2,1$\\
$0<\rho<\lc$ & one & simple\\
$\rho=\lc$ & two & one double, one simple\\
$\lc<\rho\leq1$ & three & all simple\\
\bottomrule
\end{tabular}
\end{center}
At the fold,
\begin{equation}\label{eq:foldnondeg}
 F_\rho(\ac,\lc)<0,\qquad F_{aa}(\ac,\lc)>0.
\end{equation}
Above it, label the zeros $\aL<\aM<\aO$. The first two lie in $(4/5,3/2)$,
while $\aO>9/5$. On $(\lc,1)$, $\aL$ and $\aO$ are strictly decreasing,
and $\aM$ is strictly increasing. As $\rho\downarrow\lc$,
\begin{equation}\label{eq:foldexp}
 \aL(\rho),\aM(\rho)=\ac\mp
 \sqrt{-\frac{2F_\rho(\ac,\lc)}{F_{aa}(\ac,\lc)}}\sqrt{\rho-\lc}
 +O(\rho-\lc).
\end{equation}
\end{theorem}
\begin{proof}
Only the interval $(4/5,3/2)$ can contain zeros other than the unique outer
one, by \cref{cor:separator,thm:outer}. Let $S$ be the set of parameters
for which $F(a,\rho)<0$ somewhere in that interval.
\Cref{lem:rhocross} makes $S$ an upper interval. \Cref{lem:lowrho} makes it
empty for $0<\rho\leq9/10$, with a positive uniform margin at $9/10$.
The negative anchor at $a=1,\rho=1$ makes it nonempty near one.
Thus $S=(\lc,1]$ with $9/10<\lc<1$.

The positive endpoint values at $4/5,3/2$, uniform continuity, and compactness
show that at $\lc$ the nonnegative function on this interval attains zero
at an interior point $\ac$. Its zero has even multiplicity at least two.
The outer simple zero and \cref{prop:three} force multiplicity exactly two
and uniqueness of $\ac$. This proves $F_{aa}>0$; \cref{lem:rhocross} gives
$F_\rho<0$.

Below $\lc$ a hypothetical left zero would be a zero of a nonnegative
function. Its strictly negative $\rho$-derivative would create negativity
at a slightly larger parameter still below $\lc$, a contradiction.
Above $\lc$, a negative value between positive endpoints gives two
sign-changing left zeros. The global multiplicity bound makes both
simple and excludes all additional zeros. At their crossings $F_a$ is
negative and positive, respectively; the common inequality $F_\rho<0$
gives the stated velocity directions. The fold expansion follows by
substituting $a=\ac+u\sqrt{\rho-\lc}$ into the analytic Taylor series.
Its limiting quadratic has two distinct roots, which continue analytically
in $\sqrt{\rho-\lc}$. Finally, at zero parameter the elementary formula
has exactly the two displayed zeros with the stated multiplicities.
\end{proof}

\Cref{thm:phase} proves precisely the global exclusion left open in the
companion's Conjecture 12.1. The local fold argument alone would not have
proved it; \cref{thm:outer} is the missing global input.

\section{Integer-gap rigidity and the exact unit-span maximum}
\label{sec:rigidity}
A classical shift identity now yields a stronger, less obvious geometric
conclusion.

\begin{theorem}[Integer-separated zeros, all odd indices]\label{thm:integergap}
Let $n\geq3$ be odd and $0<\rho\leq1$. If two positive zeros of
$U_n(\cdot,\rho)$ differ by a positive integer, they are exactly $1$ and $2$.
In particular, no positive integer gap of size at least two is possible.
\end{theorem}
\begin{proof}
For odd $n$, $g_n(x)$ is nonnegative on $(0,e^n)$, with its only zero there
at $1$. It is strictly negative on $(e^n,\infty)$. Hence every positive
zero lies below $e^n$, and $e^n$ itself is not a zero for $\rho>0$.
Iterating the shift equation gives
\[
 \rho^q U_n(a+q,\rho)=U_n(a,\rho)-
             \sum_{j=0}^{q-1}\rho^j g_n(a+j),\qquad q\in\N.
\]
If both endpoints are zeros, all arguments in the finite sum lie below
$e^n$, so every summand is nonnegative. Vanishing forces every one to be
zero. This is possible only for $q=1,a=1$.
\end{proof}

\begin{theorem}[Sharp unit-span bound and resonance]\label{thm:span}
For $\lc<\rho\leq1$, let
\[
 D(\rho)=\aO(\rho)-\aL(\rho).
\]
Then
\begin{equation}\label{eq:span}
 D(\rho)\leq1.
\end{equation}
There is a unique $\rs\in(\lc,1)$ at which equality holds. At this parameter,
\begin{equation}\label{eq:lockedzeros}
 \aL(\rs)=1,\qquad \aO(\rs)=2.
\end{equation}
The bound extends by continuity to the fold, where it is strict.
\end{theorem}
\begin{proof}
Since $\aL>4/5$, one has $\aL+1>9/5>\aM$.
The shift equation and $f(\aL)\geq0$ give
\[
 F(\aL+1,\rho)=-\frac{f(\aL)}{\rho}\leq0.
\]
Between $\aM$ and $\aO$ the function is strictly positive, and after
$\aO$ it is strictly negative. Therefore $\aL+1\geq\aO$, proving
\eqref{eq:span}. Equality holds exactly when $f(\aL)=0$, hence exactly
when $\aL=1$.

At $a=1$ the first nonzero coefficient of $F(1,\rho)$ is
$\rho f(2)>0$, whereas $F(1,1)<0$. \Cref{lem:rhocross} gives a unique
$\rs\in(0,1)$ with $F(1,\rs)=0$. The shift equation then gives
$F(2,\rs)=0$. The phase diagram excludes $\rs<\lc$.
Nor can $\rs=\lc$: if $1$ were the double fold zero, differentiating
\eqref{eq:shift} and using $f'(1)=0$ would also make $2$ a multiple zero,
contradicting the multiplicity bound of three. Thus $\rs>\lc$.
The zero at $2$ is the outer one. If the zero at $1$ were the middle one,
then the span from the left zero to $2$ would exceed one. Hence it is the
left zero. Uniqueness of the parameter zero proves uniqueness of equality.
At the fold, equality would again force a double zero at $1$ and a
multiple zero at $2$, so the limiting span is strictly smaller than one.
\end{proof}

\begin{theorem}[Common velocity and nondegenerate span maximum]\label{thm:curvature}
At the resonance, the two locked branches have the same nonzero velocity:
\begin{equation}\label{eq:velocitylock}
 v:=\aL'(\rs)=\aO'(\rs)
   =-\frac{F_\rho(1,\rs)}{F_a(1,\rs)}<0.
\end{equation}
Moreover,
\begin{equation}\label{eq:curvature}
 D''(\rs)=\frac{6v^2}{F_a(1,\rs)}<0.
\end{equation}
In particular, the unique global maximum $D(\rs)=1$ is nondegenerate, and
\begin{equation}\label{eq:spantaylor}
 D(\rho)=1+\frac{3v^2}{F_a(1,\rs)}(\rho-\rs)^2
                    +O((\rho-\rs)^3).
\end{equation}
\end{theorem}
\begin{proof}
At a simultaneous zero at $1,2$, differentiation of the shift equation
in $a$ and $\rho$ gives
\[
 F_a(2,\rs)=\frac{F_a(1,\rs)}{\rs},\qquad
 F_\rho(2,\rs)=\frac{F_\rho(1,\rs)}{\rs}.
\]
Here $f'(1)=0$ and $F(1,\rs)=0$ remove the extra terms. Implicit
differentiation gives the common velocity.

Put $\delta(\rho)=1-D(\rho)=\aL+1-\aO$. At the resonance,
$\delta=\delta'=0$. Taylor expansion at the outer simple zero yields
\[
 F(\aL+1,\rho)=F_a(2,\rs)\delta+O((\rho-\rs)^3),
\]
where analyticity and $\delta=O((\rho-\rs)^2)$ justify the remainder.
On the other hand,
$f(1+u)=3u^2+O(u^3)$ and the shift identity give
\[
 F(\aL+1,\rho)=-\frac{3v^2}{\rs}(\rho-\rs)^2
                         +O((\rho-\rs)^3).
\]
Use $\rs F_a(2,\rs)=F_a(1,\rs)$ and compare coefficients. Since the
left zero has $F_a<0$ and $v\ne0$, the curvature is strictly negative.
\end{proof}

The maximum theorem does not, by itself, prove that the span is strictly
increasing before $\rs$ and strictly decreasing afterwards. Additional
smaller local extrema are not excluded by a unique global maximum. This
stronger shape question is kept separate in \cref{sec:questions}.

\section{Polylogarithmic identities and strict inequalities}
\label{sec:polylog}
For clarity, angle brackets denote differentiation in order:
\[
 \ordli{s}{j}(\rho)=\left.\partial_z^j\Li_z(\rho)\right|_{z=s}.
\]
They never denote repeated differentiation in the argument $\rho$.
For $0<\rho<1$ these derivatives follow from the absolutely convergent
power series, and commute with argument differentiation.

\subsection{A unique order-derivative resonance}
Define
\begin{align}
 B(\rho)&=\ordli{2}{3}(\rho)+3\ordli{2}{2}(\rho),\label{eq:B}\\
 C(\rho)&=\ordli{1}{3}(\rho)+3\ordli{1}{2}(\rho),\label{eq:C}\\
 T(\rho)&=2\ordli{3}{3}(\rho)+9\ordli{3}{2}(\rho)
                                      +6\ordli{3}{1}(\rho).\label{eq:T}
\end{align}
Direct comparison of coefficients gives
\begin{equation}\label{eq:BTrelation}
 F(1,\rho)=\frac{B(\rho)}\rho,\qquad
 F_a(1,\rho)=-\frac{T(\rho)}\rho,\qquad
 B'(\rho)=\frac{C(\rho)}\rho.
\end{equation}

\begin{corollary}[Root-locked polylogarithmic identities]\label{cor:polyidentities}
The equation
\begin{equation}\label{eq:resonanceidentity}
 \boxed{\ordli{2}{3}(\rs)=-3\ordli{2}{2}(\rs)}
\end{equation}
has exactly one solution $\rs\in(0,1)$. It is precisely the unit-span
resonance of \cref{thm:span}, and $\lc<\rs<1$.
At this parameter,
\begin{align}
 \boxed{\aL'(\rs)=\aO'(\rs)=\frac{C(\rs)}{\rs T(\rs)}} ,
 \label{eq:polyvelocity}\\
 \boxed{D''(\rs)=-\frac{6C(\rs)^2}{\rs T(\rs)^3}} .
 \label{eq:polycurvature}
\end{align}
One has $C(\rs)<0$ and $T(\rs)>0$.
\end{corollary}
\begin{proof}
The unique parameter zero of $F(1,\rho)$ was proved in \cref{thm:span};
\eqref{eq:BTrelation} identifies it with \eqref{eq:resonanceidentity}.
At the zero, $F_\rho(1,\rs)=B'(\rs)/\rs=C(\rs)/\rs^2<0$.
Since $F_a(1,\rs)<0$, the signs of $C,T$ follow.
Substitution into \cref{thm:curvature} gives the two formulas.
\end{proof}

These are exact identities at a uniquely characterized interior parameter,
not a claim that $\rs$ is algebraic or that its value reduces to a previously
named constant. The uniqueness, simultaneous integer zeros, and curvature
formula carry the additional content; merely defining a constant as a zero
of $B$ would not establish those facts.

There is also a parameter-interval inequality without an implicitly defined
special value. Since $f'(1)=0$, differentiating the shift equation gives
$F_a(2,\rho)=F_a(1,\rho)/\rho$. \Cref{lem:strip} therefore proves
\begin{equation}\label{eq:Tpositive}
 \boxed{2\ordli{3}{3}(\rho)+9\ordli{3}{2}(\rho)+6\ordli{3}{1}(\rho)>0
 \quad(9/10\leq\rho\leq1).}
\end{equation}
All order derivatives in this display converge at the endpoint $\rho=1$.

\subsection{An all-index tail-moment inequality}
For an integer $M\geq1$, set
\begin{equation}
 L_{s,j}^{[M]}(\rho)=\sum_{r=M}^\infty\frac{\rho^r\log^jr}{r^s}
 =(-1)^j\ordli{s}{j}(\rho)-
                  \sum_{r=1}^{M-1}\frac{\rho^r\log^jr}{r^s}.
\end{equation}
At $s=2,3$ all these sums converge even when $\rho=1$.

\begin{corollary}[Strict polylogarithmic tail inequalities]\label{cor:tailineq}
Let $n\geq2$, $0<\rho\leq1$, and let $M$ be an integer with
$M\geq e^{\xi_n}$. Then
\begin{align}\label{eq:tailineq}
 &2L_{3,n}^{[M]}-3nL_{3,n-1}^{[M]}+n(n-1)L_{3,n-2}^{[M]}\notag\\
 &\hspace{25mm}+e^{-n}\bigl(nL_{2,n-1}^{[M]}-L_{2,n}^{[M]}\bigr)<0.
\end{align}
The simpler sufficient condition $M\geq e^{(n-1)/2}$ avoids the implicit
cutoff. In particular, for every integer $M\geq3$,
\begin{equation}\label{eq:cubic-tailineq}
 2L_{3,3}^{[M]}-9L_{3,2}^{[M]}+6L_{3,1}^{[M]}
       +e^{-3}(3L_{2,2}^{[M]}-L_{2,3}^{[M]})<0.
\end{equation}
\end{corollary}
\begin{proof}
At $a=M$, multiply the differential inequality of \cref{thm:sharp} by
$\rho^M>0$ and change summation variable from $m$ to $r=M+m$.
Differentiation of \eqref{eq:gn} supplies the three order-three moments.
The endpoint equality of the pointwise inequality, if it occurs, is made
strict by later positive-weight summands. The cubic sufficient condition
follows already from $e<3$ and \cref{thm:barrier}.
\end{proof}

These inequalities furnish sign tests for combinations of polylogarithmic
order derivatives and finite logarithmic sums. They are different in kind
from rational reductions of cyclotomic multiple polylogarithms: the relevant
structure here is a positive real deformation and a first-order barrier.

\section{Proof-bearing arithmetic and independent diagnostics}
\label{sec:certificate}
The finite parts of the proof are intentionally small enough to reproduce
without a computer algebra system. The verifier uses only Python's standard
library. This section specifies the mathematics behind it so that the
certificate is not an unexplained software output.

\subsection{Integer interval arithmetic and logarithms}
An interval is stored as $[l/S,u/S]$ with integers $l\leq u$ and
$S=10^{60}$. Each operation rounds down its lower endpoint and rounds up its
upper endpoint. Multiplication takes the minimum and maximum of the four
endpoint products. Reciprocal intervals reject a zero-containing input.
Thus induction on the computation tree proves inclusion for every result.
No floating-point comparison decides a proof-bearing sign.

For a positive rational $x$, range reduction writes $x=2^k v$ with
$1\leq v\leq2$. With $z=(v-1)/(v+1)\in[0,1/3]$,
\begin{equation}\label{eq:logseries}
 \log v=2\sum_{j=0}^{J-1}\frac{z^{2j+1}}{2j+1}+R_J,
 \qquad
 0\leq R_J\leq
 \frac{2(1/3)^{2J+1}}{(2J+1)(1-1/9)}.
\end{equation}
The code uses $J=68$, computes the finite sum with outward interval
operations, and adds the explicit rational tail. The same formula gives
$\log2$, after which $\log x=k\log2+\log v$. On an input interval,
monotonicity of logarithm supplies the endpoint enclosure. This is valid
for intervals as well as point inputs; it is why the derivative strip
covers every $a$, not only the cell endpoints.

\subsection{Euler--Maclaurin anchors with a stated remainder}
Let $P_0(y)=3y^2-y^3$, and recursively define
\begin{equation}\label{eq:polyrec}
 P_{r+1}(y)=P_r'(y)-(r+2)P_r(y).
\end{equation}
Then $f^{(r)}(x)=x^{-r-2}P_r(\log x)$ exactly.
For $b=a+N$ and a positive integer $p$, Euler--Maclaurin gives
\begin{align}\label{eq:EM}
 F(a,1)={}&\sum_{m=0}^{N-1}f(a+m)-\frac{\log^3b}{b}
 +\frac{f(b)}2\notag\\
 &-\sum_{j=1}^p\frac{B_{2j}}{(2j)!}f^{(2j-1)}(b)+R,
\end{align}
with
\begin{equation}\label{eq:EMerror}
 |R|\leq\frac{|B_{2p}|}{(2p)!}
                 \int_b^\infty |f^{(2p)}(x)|\dd x.
\end{equation}
The integral term in \eqref{eq:EM} is exact because
$\partial_x(\log^3x/x)=f(x)$ and $\log^3x/x\to0$.
The standard periodic-Bernoulli remainder yields \eqref{eq:EMerror};
its bound $|\widetilde B_{2p}(x)|\leq|B_{2p}|$ also follows immediately
from the absolutely convergent cosine Fourier series for $p\geq1$
\cite{dlmf-em}. All derivatives at infinity vanish here.

If $P_{2p}(y)=\sum_j c_jy^j$ and $b\geq1$, the absolute integral is at
most $\sum_j|c_j|I_{2p+2,j}(b)$, where repeated integration by parts gives
\begin{equation}\label{eq:logintegral}
 I_{d,j}(b)=\int_b^\infty\frac{\log^jx}{x^d}\dd x
 =b^{1-d}\sum_{h=0}^j
   \frac{j!}{(j-h)!(d-1)^{h+1}}\log^{j-h}b.
\end{equation}
This is a finite interval expression. The four anchors use $N=32,p=6$;
Bernoulli numbers are constructed as exact rationals from their defining
recurrence. The certificate includes the actual remainder enclosure for
each anchor, not only the final sign.

\subsection{What the machine certificate does and does not prove}
The verifier reconstructs four endpoint enclosures, the elementary rational
low-$\rho$ margin, all $128$ derivative-strip bounds, the common tail bound,
and the exponential rational estimates used in the proof. A separate test
file checks the interval primitives against exact rational arithmetic and
checks the differentiation polynomials.
The infinite-parameter consequences are the theorems above; they are not
inferred by extrapolating the finite cell count.

The file \texttt{certificates/exact.json} is generated evidence, not a
trusted oracle. Running the verifier recomputes it from the formulas.
The mathematical proof still requires the stated convergence, variation,
and barrier arguments. Conversely, none of those arguments excuses a
failed finite sign check.

\subsection{Independent numerical diagnostics and a differentiation trap}
The optional \texttt{numerical\_diagnostics.py} uses the defining Laplace
integral, with $\mu=-\log\rho$, rather than the finite Euler--Maclaurin
anchor procedure. It gives the following \emph{non-certified} approximations:
\begin{align*}
 \lc&\approx0.9999904668377386342470714889616,\\
 \ac&\approx1.049753050880255458339760685838,\\
 \rs&\approx0.9999923567014181627908364969005,\\
 \aM(\rs)&\approx1.111482903903473928636617026947.
\end{align*}
At the resonance, the common branch velocity is approximately
$-12527.90098335$ and $D''(\rs)\approx-4.26341756559\times10^9$.
These decimals are not isolating intervals and are not used in any proof.
Their large scales explain why a coarse grid in $\rho$ can miss the fold
and the narrow maximum.

A default tiny-step numerical differentiation of the polylogarithm in its
order gave a residual of about $1.28\times10^{-17}$ in the resonance check
in this environment despite a 40-decimal working precision. That is not
evidence against the identity and is not a 40-digit verification.
The diagnostics therefore also differentiate by Cauchy/Fourier coefficient
extraction on a radius-$1/4$ circle, using $128$ shifted angular nodes.
Its residual and the independent Laplace residual are recorded separately.
This comparison is still numerical, not an interval quadrature certificate.
The script preserves the naive residual so that loss of effective accuracy
is visible rather than silently discarded.

\section{Audit notes and integration corrections}
\label{sec:audit}
\subsection{An already reported notation error remains in the inspected source}
Near the end of the classical half-unit proof in
\texttt{09-zero-geometry.tex}, the sentence saying that the root of $P_1$
is $-\gamma$ uses the wrong polynomial label. In that chapter,
$P_k(t)=\prod_{j=1}^k(1+t/j)$, so $P_1(t)=1+t$ has root $-1$.
The intended polynomial is $Q_1(x)=x+\gamma$. The companion already
identified this issue; we confirm that it remains in the inspected blob
rather than presenting it as a newly discovered mathematical flaw.
The proposed correction is $P_1\mapsto Q_1$ in that sentence.

The same proof begins with a duplicated phrase, ``The endpoint signs in
the endpoint signs and the Laplace representation above''. The package
proposes a local prose correction. Neither edit changes a theorem.
Both replacements are explicit in \texttt{integration/corrections.json};
the helper emits a diff by default and checks the baseline blob before
an explicitly requested local application.

\subsection{Update the conjecture status, but not its history}
The companion's localized-fold theorem is correctly restricted. Its
Conjecture 12.1 should now be marked proved by \cref{thm:outer,thm:phase},
not retroactively described as following from the earlier local argument.
The missing input was a global outer exclusion. The integration fragment
preserves that dependency and uses new label prefixes to avoid collisions.

\subsection{Normalization and endpoints must remain explicit}
The four anchors in this article enclose $F$ itself. Some manuscript tables
instead enclose $a^2F$ at derivative order one. For example the earlier
normalized value at $a=3/2$ is approximately $0.05603756$, whereas the
unnormalized value here is approximately $0.02490558$; these are consistent,
not a discrepancy. Similarly, zero counts with multiplicity differ from
counts of distinct zeros at $\rho=0$ and at the fold.

The all-index sharp barrier has equality in its \emph{differential}
expression at its pointwise cutoff. Strictness at the cutoff for a
geometric sum requires $\rho>0$. We state this explicitly rather than
extending a strict inequality to a boundary where it is false.

\section{Further research questions and proposed conjectures}
\label{sec:questions}
\paragraph{1. A stronger shape theorem for the span.}
The exact unit maximum motivates the following strengthening. It is not
proved by the present maximum theorem and is not needed by any result above.
\begin{conjecture}[Strict span unimodality]\label{conj:span}
The function $D$ is strictly increasing on $(\lc,\rs)$ and strictly
decreasing on $(\rs,1)$, with no other stationary point.
\end{conjecture}
A possible route is to differentiate the exact shifted-root equation
$F(\aL+1,\rho)=-f(\aL)/\rho$ and control the average negative outer slope.
The sign of $f(\aL)$ alone proves the maximum but does not control that
average's variation.

\paragraph{2. The optimal exterior operator.}
For each $n$, determine the least cutoff obtainable from
$g_n'+\lambda g_n<0$ after optimizing over $\lambda>0$.
\Cref{thm:sharp} optimizes the cutoff only with $\lambda=e^{-n}$ fixed.
The optimized problem can be expressed as a separation of logarithmic
slope ranges on the positive and negative parts of $g_n$.
It may substantially reduce the compact interval left for certification.

\paragraph{3. All-index outer uniqueness.}
Determine for which $n$ the largest positive zero of $U_n$ is unique in a
common exterior region for every $\rho$. The barrier supplies a uniform
one-crossing region, but not automatic positivity at its left endpoint.
The index-three argument suggests combining a few certified separators
with a sign-aware geometric-weight bound, rather than scanning the full
$(a,\rho)$ plane.

\paragraph{4. Which odd indices admit an integer resonance?}
For odd $n\geq3$, an integer-separated pair must be $1,2$ by
\cref{thm:integergap}. A necessary and sufficient condition for an interior
resonance is $\gamma_n'(1)<0$, because $U_n(1,\rho)$ starts positive
and has one coefficient sign change. A zero exactly at the endpoint
$\rho=1$ should be treated separately. Determining these signs uniformly
in $n$ would connect a rigid root geometry with oscillatory Stieltjes
coefficient data. No all-index sign pattern is asserted here.

\paragraph{5. Certify the critical parameters themselves.}
The present proof certifies the existence, uniqueness, and nondegeneracy
of the fold and the resonance, but the displayed many-digit parameter
values are diagnostics. Interval Newton or a Krawczyk enclosure, together
with a rigorously subtracted endpoint singularity, should supply narrow
rational isolating boxes. The near-endpoint variable $\mu=-\log\rho$ is
preferable to an unstructured high-order $\rho$ grid.

\paragraph{6. Higher derivative orders and formal verification.}
The unresolved uniform-saturation question asks whether the necessary
threshold $k\geq n-1$, imposed by the elementary endpoint, is sufficient
throughout the deformation. Our result does not settle it. The integer
interval engine, logarithm remainder, exact derivative-polynomial
recurrence, and rational strip bound are natural first proof-assistant
targets. Formalizing the analytic variation and compactness arguments
would then isolate a relatively small trusted analysis layer.

\paragraph{7. Sharp inequalities for order derivatives.}
The moment inequality \eqref{eq:tailineq} can be combined with optimized
operators or several shifts to produce stronger sign certificates for
polylogarithmic order derivatives. Investigate whether such inequalities
can guide integer-relation searches without confusing inequalities with
identities, and whether analogous barriers survive positive mixtures of
Lerch parameters or nonuniform sampling measures.

\section{Reproduction and repository integration}
\label{sec:reproduce}
The package is independent of the user's working checkout. From its root,
\begin{verbatim}
python code/verify_exact.py
python code/test_exact.py
python code/numerical_diagnostics.py --dps 40   # optional; requires mpmath
cd article
pdflatex -interaction=nonstopmode -halt-on-error lerch_global_phase.tex
pdflatex -interaction=nonstopmode -halt-on-error lerch_global_phase.tex
\end{verbatim}
The first command reconstructs the proof-bearing JSON. The second checks
arithmetic and polynomial regression tests. The optional third command is
not part of the proof. The final two commands compile this article.

For integration, a suggested report destination is
\path{Analysis/Polylogarithms/docs/reports/lerch-global-phase/}.
The package includes a manuscript-sized theorem fragment and instructions
for adding it at the end of the current zero-geometry chapter. Any code or fragment
should be reviewed against the actual integration commit, since the
repository may have changed after the pinned snapshot. The guarded
correction helper refuses an unexpected baseline rather than overwriting
unseen edits. No automatic remote write is part of the package.

\appendix
\section{A general existence criterion for odd-index resonances}
\label{app:oddres}
For odd $n\geq3$, the first nonzero coefficient of $U_n(1,\rho)$ is
$\rho g_n(2)>0$, because $\log2<n$. Its coefficients have precisely one
positive-to-negative sign change. The crossing argument in
\cref{lem:rhocross} applies without alteration. Therefore:
\begin{proposition}\label{prop:oddres}
For odd $n\geq3$, there is at most one $\rho\in(0,1)$ with positive zeros
at both $1$ and $2$. Such an interior parameter exists if and only if
$\gamma_n'(1)<0$. If $\gamma_n'(1)=0$, there is no interior parameter,
but the endpoint $\rho=1$ has both zeros. If $\gamma_n'(1)>0$, there is
no such parameter in $(0,1]$.
\end{proposition}
\begin{proof}
The only possible integer pair is $1,2$, and the shift equation makes its
existence equivalent to $U_n(1,\rho)=0$. The function starts positive,
has at most one interior zero with strictly negative derivative, and is
continuous at one. An interior crossing forces a strictly negative
endpoint value by the separating-power comparison, even if one attempts
to approach the endpoint from below. Conversely a negative endpoint
forces an interior zero. A zero endpoint cannot coexist with an interior
crossing, by the same strict comparison.
\end{proof}
In polylogarithmic notation the unique interior resonance, when it exists,
satisfies
\begin{equation}\label{eq:generalresonance}
 \boxed{\ordli{2}{n}(\rho)+n\ordli{2}{n-1}(\rho)=0.}
\end{equation}
For odd $n$, direct coefficient comparison identifies the left side with
$\rho U_n(1,\rho)$. Thus \eqref{eq:generalresonance} is an all-odd-index
family with an exact existence criterion, not an assertion that such a
parameter exists at every odd index. The cubic case is the one completely
resolved here, including its surrounding global zero geometry.

\section{Certificate parameters and dependency map}
\begin{center}
\begin{tabular}{ll}
\toprule
Item & Fixed value or dependency\\
\midrule
Interval denominator & $10^{60}$\\
Logarithm terms & $68$, with rational atanh tail\\
Anchor summands & $N=32$\\
Euler--Maclaurin order & $p=6$\\
Derivative cells & $128$ equal closed rational intervals\\
Derivative summands & $64$ per cell\\
Derivative parameter strip & $[9/10,1]$\\
Coarse derivative conclusion & $F_a<-1/40$\\
Fold existence & anchors, positivity, compactness, three-zero bound\\
Global exclusion & low-$\rho$ bound, derivative strip, analytic barrier\\
Unit-span maximum & global diagram, shift equation, multiplicity bound\\
Many-digit diagnostics & not used by any theorem\\
\bottomrule
\end{tabular}
\end{center}

\begin{thebibliography}{9}
\bibitem{repo}
V.~Reshetnikov, \emph{ProveIt: Analysis/Polylogarithms manuscript},
pinned commit \sha, especially
\texttt{chapters/09-zero-geometry.tex} and \texttt{chapters/10-discovery.tex}.
\url{https://github.com/VladimirReshetnikov/ProveIt/tree/cc34f73596336f2466d9754cb0f3635bd2bedade/Analysis/Polylogarithms/docs/manuscript}.

\bibitem{previous}
\emph{Lerch Endpoint Singularities and Stieltjes Zero Bifurcations},
OpenAI-assisted research continuation for Vladimir Reshetnikov's ProveIt
project, October 9, 2026, especially Theorem 8.1 and Conjecture 12.1.
Retrieved LaTeX source \texttt{lerch\_zero\_bifurcations.tex}; source hash
recorded in the accompanying provenance manifest.

\bibitem{dlmf-lerch}
NIST Digital Library of Mathematical Functions, \S25.14,
\emph{Lerch's Transcendent}, including the defining series, shift relation,
and Mellin integral. \url{https://dlmf.nist.gov/25.14}.

\bibitem{dlmf-hurwitz}
NIST Digital Library of Mathematical Functions, \S25.11,
\emph{Hurwitz Zeta Function}, including parameter differentiation and
Euler--Maclaurin representations. \url{https://dlmf.nist.gov/25.11}.

\bibitem{dlmf-gamma}
NIST Digital Library of Mathematical Functions, \S5.8,
\emph{Infinite Products for the Gamma Function}.
\url{https://dlmf.nist.gov/5.8}.

\bibitem{dlmf-em}
NIST Digital Library of Mathematical Functions, \S2.10,
\emph{Sums and Sequences}, especially the Euler--Maclaurin formula.
\url{https://dlmf.nist.gov/2.10}.

\bibitem{coffey}
M.~W.~Coffey, \emph{Series representations for the Stieltjes constants},
Rocky Mountain Journal of Mathematics \textbf{44} (2014), 443--477;
preprint arXiv:0905.1111 (2009).
\url{https://arxiv.org/abs/0905.1111}.
\end{thebibliography}
\end{document}
